
% =============================================================================
% Project : THE PHILOSOPHY OF VIABILITY
% Book    : Book I — Ontology of Viability
% Book ID : VS.B01
% Part    : Part I — Identity and Manifestation
% Chapter : Chapter 08 — Comparative Philosophical Examination
% Chapter ID : VS.B01.P01.C08
%
% File :
%   4_ASSETS/diagrams/
%   VS.B01.P01.C08.S09.SS07.P006.DG01.hume_attribution_of_personal_identity.tex
%
% Section / Chapter / Asset:
%   Diagram DG01 — Hume's Account of the Attribution of Personal Identity
%
% Structural Identifier:
%   VS.B01.P01.C08.S09.SS07.P006.DG01
%
% Asset Type : DG
%
% Purpose:
%   To represent the relations and psychological transitions
%   through which Hume explains the attribution of personal
%   identity to a succession of perceptions.
%
% Representational Limit:
%   The diagram represents Hume's explanatory account of
%   identity attribution. It does not establish that an
%   ontological identity is constituted by these psychological
%   relations or operations.
%
% Dependencies:
%   TikZ; positioning library.
%   Loaded by the project preamble.
%
% =============================================================================

\begin{figure}[htbp]
    \centering

    \begin{tikzpicture}[
        font=\small,
        node distance=0.72cm,
        process/.style={
            draw=blue!55!black,
            fill=blue!7,
            rounded corners=3pt,
            line width=0.7pt,
            text width=8.2cm,
            minimum height=0.85cm,
            align=center,
            inner sep=6pt
        },
        central/.style={
            process,
            fill=blue!14,
            font=\small\bfseries
        },
        qualification/.style={
            draw=black!45,
            fill=black!3,
            rounded corners=2pt,
            text width=11.2cm,
            align=center,
            font=\footnotesize,
            inner sep=7pt
        },
        flow/.style={
            ->,
            draw=black!65,
            line width=0.8pt
        }
    ]

    \node[process] (perceptions) {
        \textbf{Distinct and Successive Perceptions}\\
        Impressions and ideas occurring in succession
    };

    \node[process, below=of perceptions] (relations) {
        \textbf{Relations among Perceptions}\\
        Resemblance \quad Contiguity \quad Causation
    };

    \node[process, below=of relations] (memory) {
        \textbf{Memory and Custom}\\
        Recollection of past perceptions and the formation
        of habitual transitions
    };

    \node[process, below=of memory] (imagination) {
        \textbf{Easy Transition of the Imagination}\\
        Related perceptions are contemplated as though
        they belonged to one uninterrupted object
    };

    \node[central, below=of imagination] (attribution) {
        Attribution of Personal Identity
    };

    \draw[flow] (perceptions) -- (relations);
    \draw[flow] (relations) -- (memory);
    \draw[flow] (memory) -- (imagination);
    \draw[flow] (imagination) -- (attribution);

    \node[qualification, below=0.85cm of attribution] (limit) {
        \textbf{Interpretive qualification:}
        This sequence explains how identity is attributed
        in Hume's account. It does not, by itself, establish
        that the ontological identity of a person is identical
        with the relations or psychological operations shown.
    };

    \end{tikzpicture}

    \caption{Hume's account of the attribution of personal
    identity through relations among successive perceptions,
    memory, custom, and the imagination.}
    \label{fig:VS.B01.P01.C08.S09.SS07.P006.DG01}
\end{figure}

% =============================================================================
% END OF DIAGRAM DG01
% =============================================================================
